\section{Hierarchical Memory Architecture}

\label{sec:memory}

\subsection{Memory Hierarchy}

Hanzo Chat implements a three-tier memory hierarchy:

\begin{definition}[Memory Tiers]
\begin{itemize}[leftmargin=1.1em]
    \item \textbf{Session Memory} ($\mathcal{M}_S$): Full conversation history within the current session. Stored in-memory, evicted on session close. Capacity: full context window.
    \item \textbf{Project Memory} ($\mathcal{M}_P$): Facts, decisions, and summaries relevant to a specific project. Persisted in vector database. Retention: indefinite.
    \item \textbf{Global Memory} ($\mathcal{M}_G$): User preferences, frequently referenced information, and cross-project knowledge. Persisted globally. Retention: indefinite.
\end{itemize}
\end{definition}

\subsection{Memory Storage}

Each memory entry is a tuple $(k, v, \bm{e}, t, s, \rho)$ where $k$ is a unique identifier, $v$ is the content string, $\bm{e} \in \mathbb{R}^d$ is the embedding vector, $t$ is the creation timestamp, $s \in \{S, P, G\}$ is the scope, and $\rho \in [0, 1]$ is the relevance score.

Embeddings are computed using a fine-tuned E5-large model~\cite{wang2022text} ($d = 1024$) and stored in a pgvector~\cite{pgvector2023} index for efficient similarity search.

\subsection{Memory Retrieval}

Given a query embedding $\bm{e}_t$, we retrieve the top-$k$ memories from each tier:

\begin{equation}
\mathcal{R}_s = \text{Top-}k_s\left(\{m \in \mathcal{M}_s : \cos(\bm{e}_t, m.\bm{e}) > \tau_s\}\right),
\end{equation}

where $k_S = 20$, $k_P = 10$, $k_G = 5$ and $\tau_S = 0.3$, $\tau_P = 0.5$, $\tau_G = 0.6$ are tier-specific thresholds (higher thresholds for broader scopes to maintain relevance).

The combined retrieval set $\mathcal{R} = \mathcal{R}_S \cup \mathcal{R}_P \cup \mathcal{R}_G$ is re-ranked by a cross-encoder~\cite{nogueira2019passage} and truncated to fit within the context budget.

\subsection{Automatic Summarization}

When session memory exceeds a threshold (default: 75\% of context window), older conversation turns are automatically summarized:

\begin{algorithm}[H]
\caption{Adaptive Memory Summarization}
\label{alg:summarize}
\begin{algorithmic}[1]
\Require Session history $H = [h_1, \ldots, h_T]$, context budget $B$
\State $\text{tokens} \gets \text{CountTokens}(H)$
\While{$\text{tokens} > 0.75 \cdot B$}
    \State $\text{oldest} \gets H[1:\lfloor T/3 \rfloor]$ \Comment{Oldest third}
    \State $\text{summary} \gets \text{LLM.Summarize}(\text{oldest})$
    \State Extract facts $\mathcal{F} \gets \text{LLM.ExtractFacts}(\text{oldest})$
    \State Store $\mathcal{F}$ in project memory $\mathcal{M}_P$
    \State $H \gets [\text{summary}] \cup H[\lfloor T/3 \rfloor + 1 : T]$
    \State $\text{tokens} \gets \text{CountTokens}(H)$
\EndWhile
\State \Return $H$
\end{algorithmic}
\end{algorithm}

This process is transparent to the user: the full conversation appears continuous even though the underlying representation has been compressed.

\subsection{Memory Consolidation}

We periodically consolidate memories across tiers using a process inspired by hippocampal replay in neuroscience~\cite{mcclelland1995memory}:

\begin{enumerate}[leftmargin=1.4em]
    \item \textbf{Deduplication}: Merge memories with cosine similarity $> 0.95$.
    \item \textbf{Promotion}: Session memories accessed $> 3$ times are promoted to project memory.
    \item \textbf{Generalization}: Cluster project memories and extract general rules for global memory.
    \item \textbf{Decay}: Memories not accessed for 90 days have their relevance score decayed by 10\%.
\end{enumerate}

\subsection{Context Window Management}

The context window is partitioned as follows:

\begin{table}[H]
\centering
\begin{tabular}{lcc}
\toprule
\textbf{Section} & \textbf{Budget (\%)} & \textbf{Example (128K)} \\
\midrule
System prompt & 5\% & 6,400 tokens \\
Global memory & 5\% & 6,400 tokens \\
Project memory & 10\% & 12,800 tokens \\
Retrieved context & 15\% & 19,200 tokens \\
Conversation history & 50\% & 64,000 tokens \\
Tool definitions & 5\% & 6,400 tokens \\
Current query & 5\% & 6,400 tokens \\
Generation headroom & 5\% & 6,400 tokens \\
\bottomrule
\end{tabular}
\caption{Context window budget allocation for a 128K token model.}
\label{tab:context_budget}
\end{table}

